
\documentclass{../concepts/titul_zprava_p/prepareprotokolZprava} %všechny balíčky pro protokol + titulní strana

\usepackage{graphicx}
\usepackage{newfloat}
\usepackage{caption}
\usepackage{booktabs}
\usepackage{amssymb}
\usepackage{amsmath}
\usepackage{siunitx}
\usepackage{textcomp}
\usepackage{comment}
\usepackage{gensymb}
\usepackage{subcaption}
\usepackage{multirow}
\sisetup{inter-unit-product =$\cdot$}
%\sisetup{separate-uncertainty=true}

\praktikum{III} %číslo praktika (I, II nebo III)
\autor{David Němec} %jméno autora
\datum{14.\,2.\,2026} %datum měření
\cislo{18} %číslo úlohy
\nazev{Laserová dopplerovská anemometrie} %název úlohy

\DeclareFloatingEnvironment[fileext=los, placement={htbp}, name=Obrázek]{schema}

\newcommand{\nobars}{
Chybové úsečky nebyly pro~svou malou velikost pro~přehlednost vykresleny.}

\newcommand{\nohbars}{
Horizontální chybové úsečky nebyly pro~svou malou velikost pro~přehlednost vykresleny.}

\newcommand{\novbars}{
Vertikální chybové úsečky nebyly pro~svou malou velikost pro~přehlednost vykresleny.}

\newcommand{\relative}[1]{
\frac{\sigma_{#1}}{#1}
}
\newcommand{\squared}[1]{
\left( #1 \right)^2
}
\newcommand{\relsqr}[1]{
\squared{\relative{#1}}
%\left( \frac{\sigma_{#1}}{#1} \right)^2
}

\begin{document}
\renewcommand{\figurename}{Graf}
\renewcommand{\refname}{Literatura}
\renewcommand{\today}{\number\day.\,\number\month.\,\number\year}

% add another graphics path to find ZPF.pdf for title page
% now searches for images in ./ and ../concepts/titul_protokol_p/
\graphicspath{{../concepts/titul_protokol_p/}}
\maketitle

% ========================== PRACOVNÍ ÚKOL ============================
\section{Pracovní úkoly}
\begin{enumerate}
    \item Proveďte kalibraci „optické sondy anemometru“. Parametry optické sondy získáte jednak měřením vzdálenosti 
    interferenčních plošek v~průsečíku laserových paprsků metodou projekce, jednak výpočtem z~geometrie uspořádání. 
    Oba výsledky porovnejte.
    \item Připravte aparaturu k~měření rychlosti částic. Zkontrolujte chod paprsků v~detekční optice a~vymezte prostorovou 
    dírkovou clonu.
    \item Na základě průběhu měřeného signálu optimalizujte dopplerovský signál na~proudění vody v~kyvetě.
    \item Změřte frekvence dopplerovského signálu na~souboru 60 -- 80 částic. Převeďte hodnoty frekvence na~hodnoty rychlostí. 
    Graficky zpracujte rozložení rychlostí ve~vodě formou histogramu. Histogram fitujte funkcí normálního rozdělení a~z~ní 
    stanovte střední hodnotu rychlosti částic a~standardní odchylku nalezeného rozdělení.
    \item Diskutujte, jaký vliv na~výsledek má to, že~parametry optické sondy jsou měřeny ve~vzduchu, zatímco měření rychlostí 
    částic probíhá ve~vodě.
\end{enumerate}

% ============================= TEORIE ================================
\section{Vztahy a veličiny}

% ============= schémata asi nejsou nutná
% \begin{schema}[!htbp]
%     \centering
%     %\includegraphics[width=0.5\linewidth]{image.png}
%     \caption{Schéma PŘÍSTROJE \cite{pokyny}}
%     \label{sch:pristroj}
% \end{schema}






% ======================== VÝSLEDKY MĚŘENÍ ===========================
\section{Výsledky a zpracování měření}

% ============= podmínky měření

% ============= přístroje a jejich nejistoty
\subsection{Použité měřicí přístroje}
\begin{enumerate}
    \item PŘÍSTROJ:
    
    POZNÁMKY

\end{enumerate}

% ============= naměřené hodnoty, vypočtené veličiny, grafy, výsledky
\subsection{Naměřené hodnoty a výpočty}

\begin{figure}
    \centering
    %\includegraphics[width=0.5\linewidth]{graf.png}
    \caption{NAZEVGRAFU}
    \label{fig:graf}
\end{figure}


% ============================ DISKUZE a ZÁVĚR ================================
\section{Diskuze výsledků a závěr}
% ======================== diskuze

% ======================== závěr


% =========================== LITERATURA =============================
\begin{thebibliography}{9}

\bibitem{studijni} Kolektiv ZFP KVOF MFF UK: \textit{NÁZEV} [online]. 
[cit. \today], dostupné z~URL

\bibitem{pokyny} Kolektiv ZFP KVOF MFF UK: \textit{NÁZEV} [online]. [cit. \today], dostupné
z~URL

\bibitem{vlnovedelky} Kolektiv ZFP KVOF MFF UK: \textit{Úloha 16 - studijní text} 
[online]. [cit. \today], dostupné
z~https://physics.mff.cuni.cz/vyuka/zfp/\_media/zadani/texty/txt\_316.pdf

\bibitem{englich} J. Englich: \textit{Úvod do praktické fyziky I: Zpracování výsledků měření}. 
1. vyd. Praha: Matfyzpress, 2006

\bibitem{mikulcak} Mikulčák a kolektiv: \textit{Matematické, fyzikální a~chemické tabulky pro~střední 
školy: NÁZEVTABULKY}, Praha: Státní pedagogické nakladatelství, n. p., 1988

\end{thebibliography}
\end{document}